\pdfminorversion=6
\ifx\fhformat\undefined\def\fhformat{}\fi
\documentclass[\fhformat,10pt]{style/fh-ooe}

\begin{document}

% 1: title page generated automatically by fh-ooe.cls

\begin{frame}{Leitfrage des Vortrags}
\Large
\begin{center}
\textbf{Wie wird aus einem persönlichen KI-Coding-Assistenten\\
\color{fh-ooe-red}ein gemeinsamer, kontrollierbarer Arbeitsraum für Lehre?}
\end{center}
\vspace{0.5em}
\normalsize
\begin{columns}[T]
\column{0.32\textwidth}
\begin{block}{Nicht der Fokus}
Noch einmal erklären, dass KI Code erzeugen kann.
\end{block}
\column{0.32\textwidth}
\begin{block}{Der Fokus}
Engineering-Prozesse explizit machen: Kontext, Rollen, Aktionen, Review.
\end{block}
\column{0.32\textwidth}
\begin{block}{Die Erweiterung}
Mehrere Personen arbeiten am selben oder an getrennten KI-Arbeitskontexten.
\end{block}
\end{columns}
\end{frame}

\begin{frame}{Vom Chatfenster zum Engineering-System}
\begin{center}
\begin{tikzpicture}[node distance=1.25cm,>=Latex, every node/.style={font=\small}]
\node[draw,rounded corners,fill=dachslight,minimum width=3.2cm,minimum height=1cm] (prompt) {Prompt / Antwort};
\node[draw,rounded corners,fill=dachslight,minimum width=3.2cm,minimum height=1cm,right=2.0cm of prompt] (workspace) {Engineering-Workspace};
\draw[-{Latex[length=3mm]},very thick,fh-ooe-red] (prompt) -- node[above]{explizit machen} (workspace);
\node[below=0.45cm of prompt,align=center,text=dachsgray] {individuell\\kurzlebig\\UI-zentriert};
\node[below=0.45cm of workspace,align=center] {Threads und Turns\\Tools und Ausführung\\Historie und Rollen\\mehrere Frontends};
\end{tikzpicture}
\end{center}
\vspace{0.8em}
\begin{alertblock}{These}
Für die Lehre ist nicht nur die Qualität der KI-Antwort interessant, sondern die \textbf{Struktur des gemeinsamen Arbeitsprozesses}.
\end{alertblock}
\end{frame}

\begin{frame}{Systemarchitektur: dünne Schichten, klare Autorität}
\centering
\begin{tikzpicture}[
  box/.style={draw,rounded corners,minimum width=3.3cm,minimum height=0.9cm,align=center,fill=dachslight},
  smallbox/.style={draw,rounded corners,minimum width=2.65cm,minimum height=0.8cm,align=center,fill=white},
  arr/.style={-{Latex[length=2.5mm]},thick},
  node distance=0.65cm and 0.8cm,
  font=\small]
\node[smallbox] (native) {CodexUI\\Qt 6 / nativ};
\node[smallbox,right=of native] (web) {CodexWebUI\\Browser};
\node[smallbox,right=of web] (other) {weitere Clients\\SNode.C / TS};
\node[box,below=0.85cm of web,minimum width=7.3cm] (bridge) {AISuite / \texttt{codex-bridge}\\Transport, Routing, Controller/Observer};
\node[box,below=0.75cm of bridge,minimum width=5.0cm] (server) {OpenAI Codex app-server\\semantische Autorität};
\node[box,below=0.75cm of server,minimum width=3.3cm,fill=fh-ooe-red!8] (codex) {Codex};
\draw[arr] (native) -- (bridge);
\draw[arr] (web) -- (bridge);
\draw[arr] (other) -- (bridge);
\draw[arr,<->] (bridge) -- node[right]{natives JSON-RPC} (server);
\draw[arr,<->] (server) -- (codex);
\node[draw,dashed,rounded corners,fit=(native)(web)(other)(bridge),inner sep=0.28cm,label={[font=\scriptsize]left:SNode.C-basierte Integrationsseite}] {};
\end{tikzpicture}
\smallsource{AISuite: stateless multi-client bridge; app-server bleibt Autorität für Threads, Turns und Items.}
\end{frame}

\begin{frame}{Warum diese Architektur Kollaboration ermöglicht}
\begin{columns}[T]
\column{0.48\textwidth}
\begin{block}{Ein app-server-Kontext}
\begin{itemize}
  \item genau eine Provider-Session pro Bridge
  \item mehrere Frontend-Verbindungen
  \item Notifications werden an alle verbundenen Frontends verteilt
  \item keine zweite semantische Kopie im Backend
\end{itemize}
\end{block}
\column{0.48\textwidth}
\begin{block}{Explizite Rollen}
\begin{itemize}
  \item ein \key{Controller}
  \item weitere Clients als \key{Observer}
  \item Kontrollübergabe explizit
  \item kein stilles ``irgendwer übernimmt'' beim Disconnect
\end{itemize}
\end{block}
\end{columns}
\vspace{0.5em}
\begin{alertblock}{Für die Lehre interessant}
\textbf{Beobachten, erklären, übernehmen, zurückgeben} werden zu technisch abbildbaren Rollenwechseln -- nicht nur zu sozialen Absprachen neben einem Einzelplatz-Tool.
\end{alertblock}
\end{frame}

\begin{frame}{CodexUI und CodexWebUI: zwei Zugänge, ein Verhalten}
\begin{center}
\begin{tikzpicture}[font=\small,>=Latex]
\node[draw,rounded corners,fill=dachslight,minimum width=3.7cm,minimum height=1.1cm,align=center] (qt) at (-3.4,0) {CodexUI\\Native Qt 6};
\node[draw,rounded corners,fill=dachslight,minimum width=3.7cm,minimum height=1.1cm,align=center] (web) at (3.4,0) {CodexWebUI\\Browser};
\node[draw,rounded corners,fill=fh-ooe-red!8,minimum width=4.4cm,minimum height=1.0cm,align=center] (bridge) at (0,-2) {codex-bridge};
\draw[<->,thick] (qt) -- (bridge);
\draw[<->,thick] (web) -- (bridge);
\node[align=center,text=dachsgray] at (0,0.35) {gleiche app-server-Semantik\\kein zusätzlicher Backend-State};
\end{tikzpicture}
\end{center}
\begin{itemize}
  \item Lehrender kann ohne lokale Qt-Installation über den Browser dazukommen.
  \item Studierende können nativ arbeiten, ohne einen eigenen Web-Backend-Stack betreiben zu müssen.
  \item Native und Web sind \textbf{Präsentationsformen desselben Codex-Arbeitsmodells}.
\end{itemize}
\end{frame}

\begin{frame}{Kollaboration als didaktischer Baustein}
\centering
\begin{tikzpicture}[
 person/.style={draw,circle,minimum size=0.95cm,fill=white},
 work/.style={draw,rounded corners,minimum width=3.5cm,minimum height=1.0cm,fill=dachslight,align=center},
 >=Latex,font=\small]
\node[person] (teacher) at (-4,1.6) {L};
\node[person] (s1) at (-4,0) {S1};
\node[person] (s2) at (-4,-1.6) {S2};
\node[work] (ws) at (0,0) {gemeinsamer\\Codex-Arbeitskontext};
\node[work,fill=fh-ooe-red!8] (ai) at (4,0) {app-server\\+ Codex};
\draw[<->,thick] (teacher) -- node[above,sloped]{beobachten / steuern} (ws);
\draw[<->,thick] (s1) -- (ws);
\draw[<->,thick] (s2) -- (ws);
\draw[<->,very thick,fh-ooe-red] (ws) -- (ai);
\end{tikzpicture}
\vspace{0.6em}
\begin{block}{Entscheidend}
Der gemeinsame Gegenstand ist nicht nur der Quellcode im Repository, sondern auch der \textbf{laufende KI-gestützte Engineering-Kontext}: was analysiert wurde, welche Aktion läuft und wer gerade steuern darf.
\end{block}
\end{frame}

\begin{frame}{Szenario A: Gemeinsames Debugging im Labor}
\begin{columns}[T]
\column{0.54\textwidth}
\centering
\begin{tikzpicture}[font=\scriptsize,>=Latex]
\node[draw,circle] (l) at (0,2.0) {L};
\node[draw,circle] (s1) at (-2.1,0.8) {S1};
\node[draw,circle] (s2) at (0,0.8) {S2};
\node[draw,circle] (s3) at (2.1,0.8) {S3};
\node[draw,rounded corners,fill=dachslight,minimum width=4.0cm,minimum height=0.9cm,align=center] (b) at (0,-0.6) {ein codex-bridge\\ein app-server};
\node[draw,rounded corners,fill=fh-ooe-red!8,minimum width=2.6cm,minimum height=0.8cm] (c) at (0,-2.0) {Codex};
\foreach \x in {l,s1,s2,s3} \draw[<->] (\x) -- (b);
\draw[<->,thick] (b) -- (c);
\end{tikzpicture}
\column{0.44\textwidth}
\begin{block}{Ablauf}
\begin{enumerate}
  \item Student erklärt Hypothese.
  \item Gruppe beobachtet Analyse/Tools.
  \item Controller wird gezielt übergeben.
  \item anderer Student testet Gegenhypothese.
  \item Lehrender greift nur bei Bedarf ein.
\end{enumerate}
\end{block}
\end{columns}
\begin{alertblock}{Mehrwert}
Eine KI-Session wird zum \textbf{gemeinsamen Whiteboard für Software Engineering} -- aber mit realem Code, Tools und Tests.
\end{alertblock}
\end{frame}

\begin{frame}{Szenario B: Individuelle Arbeitsräume, gezielte Betreuung}
\centering
\begin{tikzpicture}[font=\scriptsize,>=Latex,
  pod/.style={draw,rounded corners,fill=dachslight,minimum width=2.4cm,minimum height=1.2cm,align=center}]
\node[draw,circle,fill=white] (l) at (0,2.6) {L};
\node[pod] (p1) at (-4,0) {S1\\bridge + app-server};
\node[pod] (p2) at (0,0) {S2\\bridge + app-server};
\node[pod] (p3) at (4,0) {S3\\bridge + app-server};
\node[below=0.35cm of p1] {eigener Thread-/Tool-Kontext};
\node[below=0.35cm of p2] {eigener Thread-/Tool-Kontext};
\node[below=0.35cm of p3] {eigener Thread-/Tool-Kontext};
\draw[<->,thick] (l) -- node[left]{bei Bedarf} (p1);
\draw[<->,thick,fh-ooe-red] (l) -- node[right]{aktuelle Betreuung} (p2);
\draw[<->,thick] (l) -- node[right]{bei Bedarf} (p3);
\end{tikzpicture}
\vspace{0.7em}
\begin{columns}[T]
\column{0.48\textwidth}
\key{Student bleibt Eigentümer seines Arbeitsraums}
\column{0.48\textwidth}
\key{Lehrender kann gezielt als Observer/Controller dazukommen}
\end{columns}
\vspace{0.5em}
Ideal für Praktika: Hilfe ohne Bildschirm- oder Sitzplatzwechsel; Webzugang senkt die Eintrittshürde.
\end{frame}

\begin{frame}{Szenario C: Rollenwechsel und Peer Review}
\centering
\begin{tikzpicture}[font=\small,>=Latex]
\node[draw,rounded corners,fill=dachslight,minimum width=2.6cm,minimum height=0.9cm] (a) at (-4,1.2) {Student A};
\node[draw,rounded corners,fill=dachslight,minimum width=2.6cm,minimum height=0.9cm] (b) at (-4,-1.2) {Student B};
\node[draw,rounded corners,fill=white,minimum width=3.2cm,minimum height=1.0cm,align=center] (w) at (0,0) {gemeinsamer Thread\\+ Repository};
\node[draw,rounded corners,fill=fh-ooe-red!8,minimum width=2.8cm,minimum height=0.9cm] (ai) at (4,0) {Codex};
\draw[->,thick] (a) -- node[above,sloped]{1. analysiert / implementiert} (w);
\draw[->,thick] (w) -- (ai);
\draw[->,thick] (ai) to[bend left=15] (w);
\draw[->,thick,fh-ooe-red] (b) -- node[below,sloped]{2. übernimmt / reviewt} (w);
\draw[->,dashed,thick] (a) to[bend left=25] node[left]{3. begründet} (b);
\end{tikzpicture}
\vspace{0.7em}
\begin{block}{Lehrmuster}
\textbf{Implementierer → Reviewer → Entscheider} als expliziter Rollenwechsel. Die KI bleibt gleich, aber Auftrag, Controller und Perspektive ändern sich.
\end{block}
\end{frame}

\begin{frame}{Lehrmuster statt ``Prompt Engineering''}
\begin{center}
\begin{tikzpicture}[font=\small,>=Latex,node distance=0.45cm]
\tikzset{phase/.style={draw,rounded corners,minimum width=2.15cm,minimum height=0.75cm,fill=dachslight,align=center}}
\node[phase] (p1) {Verstehen};
\node[phase,right=of p1] (p2) {Hypothese};
\node[phase,right=of p2] (p3) {Änderung};
\node[phase,right=of p3] (p4) {Test};
\node[phase,right=of p4] (p5) {Review};
\foreach \a/\b in {p1/p2,p2/p3,p3/p4,p4/p5} \draw[->,thick] (\a) -- (\b);
\end{tikzpicture}
\end{center}
\vspace{0.6em}
\begin{columns}[T]
\column{0.48\textwidth}
\begin{block}{Technisch erzwingbar / sichtbar}
\begin{itemize}
  \item eigener Turn pro Phase
  \item ``nichts ändern'' vs. Tool-Nutzung
  \item separater Review-Thread
  \item Controller-Wechsel
\end{itemize}
\end{block}
\column{0.48\textwidth}
\begin{block}{Didaktisch beobachtbar}
\begin{itemize}
  \item Hypothesenqualität
  \item Umgang mit Gegenargumenten
  \item Teststrategie
  \item begründete Annahme/Ablehnung von KI-Vorschlägen
\end{itemize}
\end{block}
\end{columns}
\end{frame}

\begin{frame}{Betreuung und Bewertung: Prozess als zusätzliches Artefakt}
\begin{alertblock}{Nicht: ``Wir loggen alles und benoten den Prompt.''}
Der Nutzen liegt darin, \textbf{Engineering-Entscheidungen in ihrem Kontext diskutierbar} zu machen.
\end{alertblock}
\vspace{0.4em}
\begin{columns}[T]
\column{0.48\textwidth}
\key{Für Betreuung}
\begin{itemize}
  \item Wo hängt der Student tatsächlich?
  \item Welche Hypothese wurde schon geprüft?
  \item Was hat die KI verändert oder ausgeführt?
  \item Kann ich zunächst nur beobachten?
\end{itemize}
\column{0.48\textwidth}
\key{Für Reflexion / Prüfung}
\begin{itemize}
  \item Warum wurde dieser Vorschlag akzeptiert?
  \item Welche Alternative wurde verworfen?
  \item Welche Tests sichern die Änderung ab?
  \item Kann der Student die Lösung ohne KI erklären?
\end{itemize}
\end{columns}
\end{frame}

\begin{frame}{Topologien für Lehrveranstaltungen}
\small
\begin{tabularx}{\textwidth}{>{\bfseries}l X X X}
\toprule
 & Arbeitskontext & Zugriff Lehrender & Geeignet für \\
\midrule
Gemeinsame Session & ein app-server / Bridge & dauerhaft Observer, gezielte Kontrolle & Demo, Debugging, Gruppenarbeit \\
\addlinespace
Individuelle Session & je Student eigener app-server / Bridge & on demand via Web oder nativ & Praktika, individuelle Betreuung \\
\addlinespace
Peer-Session & geteilter Kontext, expliziter Controller-Wechsel & optional Observer & Pair Programming, Review, Prüfungsgespräch \\
\bottomrule
\end{tabularx}
\vspace{0.8em}
\begin{block}{Technische Konsequenz}
Die Kollaborationsform wird durch \textbf{Deployment-Topologie + Rollenpolitik} definiert -- nicht durch eine neue Variante der KI.
\end{block}
\end{frame}

\begin{frame}{Was noch offen ist}
\begin{columns}[T]
\column{0.48\textwidth}
\begin{block}{Technisch}
\begin{itemize}
  \item Identität und Authentisierung bei Netzbetrieb
  \item feinere Rollen als Controller/Observer?
  \item Session Discovery im Labor
  \item Rechte auf Tools / Dateien / Shell
\end{itemize}
\end{block}
\column{0.48\textwidth}
\begin{block}{Didaktisch}
\begin{itemize}
  \item Welche Kollaborationsform verbessert Lernen messbar?
  \item Wann wird Beobachtung zu Überbetreuung?
  \item Welche Prozessdaten sind sinnvoll -- welche nicht?
  \item Wie verändert sich Prüfungsdesign?
\end{itemize}
\end{block}
\end{columns}
\vspace{0.5em}
\centering
\large \textbf{Das System ist Infrastruktur für diese Experimente -- nicht die didaktische Antwort selbst.}
\end{frame}

\begin{frame}{Live: ein Arbeitskontext, mehrere Rollen}
\Large
\begin{center}
\textbf{ca. 8--9 Minuten}
\end{center}
\normalsize
\begin{enumerate}
  \item CodexUI: bestehendes Projekt / Thread öffnen.
  \item Codex analysiert eine konkrete Aufgabe -- zunächst ohne Änderung.
  \item Änderung + Build/Test ausführen.
  \item Zweiter Client über CodexWebUI verbindet sich als Observer.
  \item Kontrolle übergeben: Review / Gegenprobe aus zweiter Perspektive.
  \item Zurückgeben und Ergebnis gemeinsam bewerten.
\end{enumerate}
\vspace{0.5em}
\begin{alertblock}{Die Demo soll genau eine Sache beweisen}
\textbf{Der KI-gestützte Engineering-Prozess kann geteilt und die aktive Rolle kontrolliert übergeben werden.}
\end{alertblock}
\end{frame}

\begin{frame}{Fazit}
\begin{center}
\Large
\textbf{Nicht noch ein KI-Frontend.}
\vspace{0.7em}

\color{fh-ooe-red}\textbf{Ein experimenteller kollaborativer Engineering-Workspace für Lehre.}
\end{center}
\vspace{0.9em}
\normalsize
\begin{itemize}
  \item app-server bleibt semantische Autorität; AISuite macht ihn netzwerkfähig und mehrbenutzerfähig.
  \item CodexUI und CodexWebUI eröffnen native und browserbasierte Zugänge zum selben Arbeitsmodell.
  \item Controller/Observer und getrennte Topologien ermöglichen neue Formen von Betreuung, Gruppenarbeit und Review.
\end{itemize}
\vfill
\centering
\textbf{Die spannende Frage für die Lehre ist nicht, was die KI kann -- sondern wie wir gemeinsam mit ihr arbeiten.}
\end{frame}

\end{document}
